% ==============================================================================
\chapter*{Abstract}
\addcontentsline{toc}{chapter}{Abstract}
% ==============================================================================

\noindent This thesis details the development of a real-time pedagogical Digital Twin for parameterized structural dynamics, bridging high-fidelity structural simulations with interactive software environments. The methodology integrates Isogeometric Analysis (IGA) with projection-based Reduced Order Models (ROM) to overcome the computational bottlenecks associated with varying geometric parameterizations. Utilizing Non-Uniform Rational B-Splines (NURBS) for exact geometric representation, dominant deformation modes are extracted via Proper Orthogonal Decomposition (POD). Furthermore, hyper-reduction strategies, specifically the Energy-Conserving Sampling and Weighting (ECSW) method, are implemented to accelerate the evaluation of parameterized internal forces. The resulting client-side web application achieves sub-16 ms solve times (60 FPS), delivering a $\mathbf{64.5\times}$ speedup and under $\mathbf{0.054\%}$ relative $L_2$ error compared to full-order models, and serves as an interactive environment for real-time parametric exploration and design.

\vspace{1.5cm}

% ==============================================================================
\chapter*{Résumé}
\addcontentsline{toc}{chapter}{Résumé}
% ==============================================================================

\noindent Ce mémoire de thèse détaille le développement d'un jumeau numérique pédagogique en temps réel pour la dynamique des structures paramétrées, reliant les simulations structurelles de haute fidélité à des environnements logiciels interactifs. La méthodologie intègre l'analyse isogéométrique (IGA) aux modèles d'ordre réduit (ROM) basés sur la projection afin de surmonter les verrous numériques associés aux paramétrisations géométriques variables. En utilisant des B-splines rationnelles non uniformes (NURBS) pour une représentation géométrique exacte, les modes de déformation dominants sont extraits par décomposition orthogonale aux valeurs propres (POD). De plus, des stratégies d'hyper-réduction, spécifiquement la méthode ECSW (Energy-Conserving Sampling and Weighting), sont implémentées pour accélérer l'évaluation des forces internes paramétrées. L'application web client finalisée atteint des temps de résolution inférieurs à 16 ms (60 FPS), offrant une accélération de $\mathbf{64.5\times}$ et une erreur relative $L_2$ inférieure à $\mathbf{0.054\%}$ par rapport aux modèles d'ordre complet, et constitue un environnement interactif pour l'exploration paramétrique et la conception en temps réel.
